% BHU PYQ -- Professional Exam Template
\documentclass[11pt,a4paper]{article}

\usepackage[a4paper,top=2.5cm,bottom=2.5cm,left=2.8cm,right=2.8cm,headheight=14pt]{geometry}
\usepackage{amsmath,amssymb}
\usepackage[T1]{fontenc}
\usepackage[utf8]{inputenc}
\usepackage{lmodern}
\usepackage{microtype}
\usepackage{fancyhdr}
\usepackage{enumitem}
\usepackage{hyperref}
\usepackage{lastpage}
\usepackage{parskip}

\hypersetup{colorlinks=true,urlcolor=black,linkcolor=black,
  pdftitle={B.Sc. (Hons.) Botany Sem-I BOB-101 -- BHU PYQ}}

\pagestyle{fancy}
\fancyhf{}
\renewcommand{\headrulewidth}{0.4pt}
\renewcommand{\footrulewidth}{0.4pt}
\fancyhead[L]{\small   \textit{Banaras Hindu University}}
\fancyhead[R]{\small   \textit{Botany Paper: BOB-101}}
\fancyfoot[C]{\small Page  \thepage\ of \pageref{LastPage}}


\newlist{parts}{enumerate}{2}
\setlist[parts,1]{label=(\alph*),leftmargin=*,itemsep=4pt,topsep=3pt}
\setlist[parts,2]{label=(\roman*),leftmargin=*,itemsep=2pt,topsep=2pt}

\newcommand{\pts}[1]{\hfill   \textit{[#1]}}


\begin{document}


\begin{center}
  {\large\bfseries BANARAS HINDU UNIVERSITY}\\[2pt]
  {
\normalsize B.Sc. (Hons.) Semester I Examination, 2016-17}\\[6pt]
  \rule{0.8\linewidth}{0.5pt}\\[6pt]
  {\large\bfseries Botany}\\[4pt]
  {\large\bfseries Paper: BOB-101 --- Cryptogams}\\[4pt]
  {
\normalsize Time: Three Hours \hfill Full Marks: 70}
\end{center}

\rule{\linewidth}{0.4pt}

\smallskip



\noindent   \textit{Note: Attempt five questions in all, selecting two questions from each Section (A and B). Question No. 9 (Section C) is compulsory. The figures in the right-hand margin indicate marks.}

\bigskip

\begin{center}
   \textbf{SECTION -- A}
\end{center}

\medskip


\noindent   \textbf{Question 1.} Describe the characteristic features of Rhodophyceae. Give an illustrated account of reproduction in   \textit{Polysiphonia}.\pts{15}

\medskip


\noindent   \textbf{Question 2.} Write short notes on any two of the following:\pts{$7.5  \times 2 = 15$}

\begin{parts}
  \item Sexual reproduction in   \textit{Sargassum}
  \item Morphology and reproduction of   \textit{Nostoc}
  \item Asexual reproduction in   \textit{Draparnaldiopsis}
\end{parts}

\medskip


\noindent   \textbf{Question 3.} With the help of labeled diagrams, give a brief account of morphology and sexual reproductive structures of   \textit{Anthoceros}.\pts{15}

\medskip


\noindent   \textbf{Question 4.} Write short notes on any two of the following:\pts{$7.5  \times 2 = 15$}

\begin{parts}
  \item Sexual reproduction in   \textit{Marchantia}
  \item Cap cell and zoospore formation in   \textit{Oedogonium}
  \item Sporangia of   \textit{Ectocarpus}
\end{parts}

\bigskip

\begin{center}
   \textbf{SECTION -- B}
\end{center}

\medskip


\noindent   \textbf{Question 5.} Give an illustrated account of various types of spore formation in   \textit{Puccinia}.\pts{15}

\medskip


\noindent   \textbf{Question 6.} Write short notes on any two of the following:\pts{$7.5  \times 2 = 15$}

\begin{parts}
  \item Reproduction in   \textit{Albugo}
  \item General characteristics of Ascomycotina
  \item Morphology of gill and basidiospore formation in   \textit{Agaricus}
\end{parts}

\medskip


\noindent   \textbf{Question 7.} With the help of suitable diagrams, give a brief account of morphology and reproduction in   \textit{Selaginella}.\pts{15}

\medskip


\noindent   \textbf{Question 8.} Write short notes on any two of the following:\pts{$7.5  \times 2 = 15$}

\begin{parts}
  \item General characteristics of Lycophyta
  \item   \textit{Alternaria}
  \item Reproduction in   \textit{Pteris}
\end{parts}

\bigskip

\begin{center}
   \textbf{SECTION -- C}
\end{center}

\medskip


\noindent   \textbf{Question 9.} Answer the following (in not more than 50 words each):\pts{$1  \times 10 = 10$}

\begin{parts}
  \item Name the reserve food material present in the Xanthophyceae.
  \item What are heterokont flagella?
  \item Name the pigments present in Cyanophyceae.
  \item Zygospore is formed in which fungus?
  \item Name the fungal group which possesses Crozier.
  \item What is the dikaryotic phase?
  \item What is the function of the receptacle in   \textit{Sargassum}?
  \item What is the function of the Vallecular canal?
  \item What are the xerophytic characteristics of   \textit{Equisetum} stem?
  \item What is the difference between an elater and a pseudoelater?
\end{parts}

\vfill
\rule{\linewidth}{0.4pt}

\begin{center}\small
\href{https://bhu-exam.vercel.app/}{Click Here}\\[4pt]\href{https://me-aryan.vercel.app/}{Click Here}\\[4pt]\href{https://quashan-me.vercel.app/}{Click Here}
\end{center}

\end{document}
